% ============================================================
% DADOS - ANEXO W - RELATÓRIO
% Arquivo: dados/dados-anexo-w-relatorio.tex
% ============================================================

\relatorioIntroducao{%
A presente sindicância foi instaurada, por determinação do Sr. \getAutoridadePosto\ \getAutoridadeNome, \getAutoridadeFuncao, por intermédio da \getPortaria, para apurar \getObjeto, conforme documentos constantes dos autos, tendo como sindicado \getSindicadoPosto\ \getSindicadoNome.
}

\relatorioDiligencias{%
Com o escopo de reunir elementos probatórios que pudessem esclarecer o fato objeto da presente sindicância, este encarregado houve por bem diligenciar conforme despacho(s) constantes dos autos, tendo sido realizadas as diligências pertinentes, incluindo-se, quando cabível, expedição e recebimento de documentos, inquirições, juntadas e demais atos necessários à elucidação dos fatos.
}

\relatorioParteExpositiva{%
Foi assegurado ao sindicado o direito ao contraditório e à ampla defesa, conforme preconizado nas Instruções Gerais para a Elaboração de Sindicância no Âmbito do Exército Brasileiro -- EB10-IG-09.001.

Da análise de todas as peças que compõem a presente sindicância, restou apurado que: [INSERIR NARRATIVA ORDENADA, COERENTE E CIRCUNSTANCIADA DOS FATOS APURADOS, COM REFERÊNCIA ÀS FOLHAS DOS AUTOS, DEPOIMENTOS, DOCUMENTOS E OUTRAS DILIGÊNCIAS].
}

\relatorioParteConclusiva{%
Em face do exposto, que dos autos consta, e conforme análise realizada na parte expositiva, verifica-se que [INSERIR CONCLUSÃO DO SINDICANTE, CONFORME UMA DAS HIPÓTESES PREVISTAS NO ANEXO W: inexistência de infração, transgressão disciplinar, indícios de crime, dano/extravio de material etc.].

Em consequência, sou de parecer que [INSERIR PROVIDÊNCIA PROPOSTA].
}

% Se não informar, usa \localData.
% \relatorioLocalData{Petrolina/PE, 31 de maio de 2026}
